\section{Belief Occupation and Coverage}
\label{supp:belief-occupation-coverage}
\label{app:coverage-details}\label{sec:coverage}\label{app:proofs-coverage}

The main paper's operational--generative value decomposition separates the
passive operating channel from the research premium.  This section develops
the passive channel through discounted belief occupation and then gives the
associated finite coverage, threshold, and belief-dynamics results.  These
are supporting results: a one-shot cost-covering region is not an optimized
Bellman research region.

Throughout, expectations and occupation rows are conditional on the displayed
initial belief.  Every horizon and state carrier is finite.  Finiteness
justifies the sum rearrangements and attained extrema below; no limiting
argument, closure-detectability condition, stationary contraction, or
path--outcome independence is used.

\subsection{Discounted gap sums and operational insertion}

\begin{definition}[Positive gap and discounted gap sum]
\label{def:frontier-gap-sum}
For a fixed candidate \(s\), define
\begin{align}
 \Delta_{s,L}(b)&:=\max\{j_s(b)-F_L(b),0\},\nonumber\\
 G_0(b;L,s)&=0,\nonumber\\
 G_{n+1}(b;L,s)&=\Delta_{s,L}(b)
 +\beta\sum_{b'}P_B(b,b')G_n(b';L,s).
 \label{eq:discounted-gap-recursion}
\end{align}
\end{definition}
Let
\[
 w_n(b,b'):=\sum_{t=0}^{n-1}\beta^t
   \Pr(B_t=b'\mid B_0=b)
\]
be discounted expected occupation, not a normalized probability.

For the supporting frozen-library process, set
\(\Delta_n^{\op,\mathrm{aux}}(s\mid b,L)
:=P_n(b,L\cup\{s\})-P_n(b,L)\), where \(P_n\) is its passive value.

\begin{proposition}[Passive gap-sum identity]
\label{thm:passive-innovation-equation}
For the supporting frozen-library process,
\begin{equation}\label{eq:passive-innovation-equation}
 \Delta_n^{\op,\mathrm{aux}}(s\mid b,L)
 =G_n(b;L,s)
 =\sum_{b'}w_n(b,b')\Delta_{s,L}(b').
\end{equation}
\end{proposition}

When the unified process has this frozen-library no-research counterfactual,
combining the occupancy identity with the operational--generative value
decomposition in the main paper gives
\begin{equation}\label{eq:strategy-innovation-equation}
 \mathcal I_n
 =\sum_{b'}w_n(b,b')\Delta_{s,L}(b')
  +\Delta_n^{\gen}.
\end{equation}
The occupancy formula shows why a zero current gap does not make a candidate
worthless if the belief process later visits states with positive gaps.

\begin{corollary}[Zero current gap, positive future operating value]
\label{cor:zero-current-gap-positive-future-value}
With a deterministic move to a gap-two belief, \(n=2\), and
\(\beta=1/2\), operational insertion value is one.
\end{corollary}

\begin{proposition}[Diminishing marginal operational innovation]
\label{prop:diminishing-operational-innovation}
If \(L\subseteq L'\), then every fixed candidate satisfies
\[
 \Delta_n^{\op}(s\mid b,L')
 \le \Delta_n^{\op}(s\mid b,L).
\]
\end{proposition}

For \cref{thm:passive-innovation-equation}, insertion replaces the frontier
pointwise by \(\max\{F_L,j_s\}\).  Subtracting \(F_L\) gives
\(\Delta_{s,L}\).  Conditional on the displayed initial belief and frozen
library, subtracting the two passive recursions and inducting on the horizon
gives \eqref{eq:discounted-gap-recursion}.  Unrolling that finite recursion
gives the discounted-occupancy representation; only finite sums are
rearranged.  The delayed-coverage specialization in
\cref{cor:zero-current-gap-positive-future-value} follows by placing unit
discounted mass on the future belief with gap two.  For
\cref{prop:diminishing-operational-innovation}, library inclusion raises the
frontier and lowers the pointwise positive gap, and nonnegative discounted
expectation preserves the order.

Persistence itself is ambiguous because it may keep the process in, or away
from, the candidate's advantage region.  The valid positive comparison is
occupation alignment on positive-gap beliefs, developed below, not
persistence in isolation.

\subsection{Coverage potential and occupation bounds}

Coverage links a candidate's operational advantage to the beliefs at which
that advantage will matter.

\begin{definition}[Certified gap and coverage potential]
\label{def:coverage-potential}
For a fixed project \(q\), library \(K\), candidate profile \(j_q\), and
operational frontier \(F_K\), define the certified gap
\[
  \Delta(q,K,b):=\max\{j_q(b)-F_K(b),0\}.
\]
Let \(\omega_t(b,b')\ge0\) be the date-\(t\) occupation weight of future
belief \(b'\) from initial belief \(b\).  For horizon \(H\), discount
\(\beta\), and candidate-survival factor \(\rho_q\), set
\begin{align}
  W_H^{\beta,\rho_q}(b,b')
    &:=\sum_{t<H}\beta^t\rho_q^t\omega_t(b,b'),
    \label{eq:discounted-coverage-occupation}\\
  \Psi_H(q,K,b)
    &:=\sum_{b'\in B}W_H^{\beta,\rho_q}(b,b')
      \Delta(q,K,b').
    \label{eq:coverage-potential}
\end{align}
\end{definition}

The potential \(\Psi_H\) is discounted exposure to beliefs at which the
candidate improves the current frontier.  The weight
\(W_H^{\beta,\rho_q}\) need not itself be a probability: it combines belief
occupation, discounting, and survival of the candidate's advantage.

\begin{theorem}[Exact finite coverage representation]
\label{thm:coverage-potential-representation}
For every finite horizon, occupation table, gap table, project, library, and
initial belief,
\begin{equation}\label{eq:coverage-representation}
  \Psi_H(q,K,b)
  =\sum_{t<H}\sum_{b'\in B}
    \beta^t\rho_q^t\omega_t(b,b')\Delta(q,K,b').
\end{equation}
The right-hand side is the date-by-date one-shot gross operational value of
the supplied candidate.
\end{theorem}

With nonnegative inputs, the same finite order argument makes
\(\Psi_H\) increasing in the gap, discount, survival, and occupation weights.
A pointwise higher frontier lowers the positive-part gap of a fixed candidate
and therefore weakly lowers its coverage potential.

\begin{corollary}[Finite occupation bounds]
\label{cor:coverage-potential-bounds}
For any nonempty \(A\subseteq B\), its attained minimum gap times its
discounted occupation is a lower bound for \(\Psi_H\).  The attained global
maximum gap times total discounted occupation is an upper bound.  If the gap
vanishes outside \(A\), the upper bound may use occupation restricted to
\(A\).
\end{corollary}

\begin{example}[Delayed coverage]
\label{ex:delayed-coverage}
On a two-belief grid, let date zero occupy the current belief and date one the
other belief.  With \(H=2\), \(\beta=1/2\), \(\rho_q=1\), current gap zero,
and future gap two,
\[
  \Delta(q,K,b_{\mathrm{current}})=0,
  \qquad
  \Psi_2(q,K,b_{\mathrm{current}})=1>0.
\]
Screening only at the current belief would discard this future advantage.
\end{example}

Fix \((q,K,b,H)\).  Substitute
\eqref{eq:discounted-coverage-occupation} into
\eqref{eq:coverage-potential}.  Since both index sets are finite,
distributivity and finite-sum commutation give
\[
\begin{aligned}
 \Psi_H(q,K,b)
 &=\sum_{b'\in B}
   \left[\sum_{t<H}\beta^t\rho_q^t\omega_t(b,b')\right]
   \Delta(q,K,b')\\
 &=\sum_{b'\in B}\sum_{t<H}
   \beta^t\rho_q^t\omega_t(b,b')\Delta(q,K,b')\\
 &=\sum_{t<H}\sum_{b'\in B}
   \beta^t\rho_q^t\omega_t(b,b')\Delta(q,K,b'),
\end{aligned}
\]
which is \eqref{eq:coverage-representation}.  This proof uses no limiting
argument.

For a nonempty region \(A\subseteq B\), write
\[
 \underline\Delta_A:=\min_{b'\in A}\Delta(q,K,b'),
 \qquad
 \overline\Delta:=\max_{b'\in B}\Delta(q,K,b').
\]
All weights and gaps are nonnegative.  Hence
\[
 \underline\Delta_A
 \sum_{b'\in A}W_H^{\beta,\rho_q}(b,b')
 \le \Psi_H(q,K,b)
 \le
 \overline\Delta
 \sum_{b'\in B}W_H^{\beta,\rho_q}(b,b').
\]
If the gap vanishes outside \(A\), the last sum can be restricted to \(A\).
The four supporting monotonicities follow term by term: raising a
nonnegative gap, discount, survival factor, or occupation weight raises every
affected summand.  If \(F_0\le F_1\), then
\(\max\{j_q-F_1,0\}\le\max\{j_q-F_0,0\}\), proving frontier saturation for a
fixed candidate.

\subsection{Monotone one-shot coverage}

For a one-step project, let \(p(b)\ge0\) scale survival or admission and let
\(\kappa(b)\) be initiation cost.

\begin{definition}[One-shot cost-covering region]
\label{def:single-gap-region}
On a finite linearly ordered belief grid with exact transition kernel \(P\),
define
\[
  G(b):=\beta p(b)\sum_{b'\in B}P(b,b')\Delta(b'),
  \qquad
  \mathcal C:=\{b\in B:\kappa(b)\le G(b)\}.
\]
The set \(\mathcal C\) is the \emph{one-shot cost-covering region}.
\end{definition}

\begin{theorem}[Monotone-gap upper-threshold theorem]
\label{thm:finite-monotone-coverage}
Assume that \(B\) is a nonempty finite linearly ordered grid; \(P\) is an
exact nonnegative row-stochastic kernel that is first-order stochastically
monotone; \(\Delta\) is nonnegative and increasing; \(p\) is nonnegative and
increasing; \(\beta\ge0\); and \(\kappa\) is antitone.  Then \(G\) is
increasing and \(\mathcal C\) is an upper set.  Consequently, either
\(\mathcal C=\varnothing\), or there is \(b^*\in B\) such that
\begin{equation}\label{eq:finite-upper-threshold}
  \mathcal C=\{b\in B:b^*\le b\}.
\end{equation}
In the latter case the region is order-connected on the finite grid.
\end{theorem}

The theorem is about the displayed one-shot cost-covering region.  We reserve
\emph{optimal research region} for the action selected by the Bellman policy.
Continuation-value differences need not inherit the order properties used
above, so the theorem does not characterize that Bellman region.

For each current belief \(b\), let
\[
 E_\Delta(b):=\sum_{b'\in B}P(b,b')\Delta(b').
\]
Here \(E_\Delta(b)\) is the conditional expected gap after one transition.
First-order stochastic monotonicity of \(P\) and monotonicity of \(\Delta\)
make it increasing.  Exact row stochasticity and \(\Delta\ge0\) make it
nonnegative.  If \(p\ge0\) is increasing, then for
\(b\le\widetilde b\),
\[
 p(b)E_\Delta(b)
 \le p(\widetilde b)E_\Delta(b)
 \le p(\widetilde b)E_\Delta(\widetilde b).
\]
Multiplication by \(\beta\ge0\) therefore makes \(G\) increasing.

Suppose \(b\in\mathcal C\) and \(b\le\widetilde b\).  Antitonicity of cost
and monotonicity of gross value yield
\[
 \kappa(\widetilde b)
 \le\kappa(b)
 \le G(b)
 \le G(\widetilde b),
\]
so \(\widetilde b\in\mathcal C\).  Thus \(\mathcal C\) is an upper set.  A
nonempty upper set in a finite linear order has a least element \(b^*\) and
equals \(\{b:b^*\le b\}\); the empty case is separate.

The cutoff directions and the exact failures of the threshold assumptions
are collected in \cref{supp:additional-comparative-statics}.

\subsection{Patience, survival, and belief dynamics}

\begin{theorem}[Finite patience--survival complementarity]
\label{thm:discount-survival-complementarity}
Let \(P\) be an exact finite row-stochastic matrix, \(g\ge0\), and
\[
  \Psi_H(\beta,\rho)
    :=\sum_{t=0}^{H-1}(\beta\rho)^tP^tg .
\]
The potential is pointwise increasing in either nonnegative \(\beta\) or
\(\rho\).  If \(0\le\beta_0\le\beta_1\) and
\(0\le\rho_0\le\rho_1\), then, componentwise,
\begin{equation}\label{eq:discount-survival-cross-difference}
  \Psi_H(\beta_1,\rho_1)+\Psi_H(\beta_0,\rho_0)
  \ge
  \Psi_H(\beta_1,\rho_0)+\Psi_H(\beta_0,\rho_1).
\end{equation}
\end{theorem}

Conditional on an initial belief, row stochasticity implies that \(P^t\) is
nonnegative, so \(P^tg\ge0\).  Monotonicity in either nonnegative scalar
follows term by term.  For the four displayed corners,
\begin{align}
 &\Psi_H(\beta_1,\rho_1)+\Psi_H(\beta_0,\rho_0)
 -\Psi_H(\beta_1,\rho_0)-\Psi_H(\beta_0,\rho_1)\notag\\
 &\qquad=
 \sum_{t<H}
 (\beta_1^t-\beta_0^t)(\rho_1^t-\rho_0^t)P^tg
 \ge0.
 \label{eq:discount-survival-factorization}
\end{align}
Indeed, the date-\(t\) coefficient factorizes exactly as
\[
 (\beta_1\rho_1)^t+(\beta_0\rho_0)^t
 -(\beta_1\rho_0)^t-(\beta_0\rho_1)^t
 =
 (\beta_1^t-\beta_0^t)(\rho_1^t-\rho_0^t).
\]
Every factor is nonnegative.  Finiteness of \(H\) justifies the displayed
termwise expansion.  No independence, inverse, infinite series, real
derivative, detectability, or contraction enters the result.

A scalar persistence parameter does not order coverage.  Consider
\[
 P(\theta)=
 \begin{pmatrix}
   \theta&1-\theta\\
   1-\theta&\theta
 \end{pmatrix},
 \qquad 0\le\theta\le1.
\]

\begin{proposition}[Opposite persistence effects]
\label{prop:no-universal-persistence-sign}
Let \(H=2\), effective discount \(\alpha=1/2\), initial belief \(0\), and
compare \(\theta_0=1/4\) with \(\theta_1=3/4\).  Exact evaluation gives
\[
\begin{array}{c@{\qquad}c@{\qquad}c}
g & \Psi_2(\theta_0) & \Psi_2(\theta_1)\\
\midrule
(1,0) & 9/8 & 11/8\\
(0,1) & 3/8 & 1/8\\
(1,1) & 3/2 & 3/2.
\end{array}
\]
Thus greater persistence can raise, lower, or leave coverage unchanged.
\end{proposition}

Persistence raises coverage when it retains occupation at an uncovered
current belief, but lowers it when it traps the process away from uncovered
beliefs elsewhere.  The positive order must therefore be stated in terms of
where occupation moves.  Define
\[
 W^P_{H,\alpha}(b,x)
 :=\sum_{t=0}^{H-1}\alpha^t(P^t)(b,x).
\]

\begin{theorem}[Advantage-region occupation alignment]
\label{thm:kernel-occupation-alignment}
Let \(P_0,P_1\) be finite row-stochastic kernels and \(g\ge0\).  If, for every
initial belief \(b\) and every \(x\) with \(g(x)>0\),
\[
  W^{P_1}_{H,\alpha}(b,x)\ge W^{P_0}_{H,\alpha}(b,x),
\]
then
\[
  \sum_{t=0}^{H-1}\alpha^tP_1^tg
  \ge
  \sum_{t=0}^{H-1}\alpha^tP_0^tg
\]
componentwise.
\end{theorem}

Fix the initial belief \(b\).  Finite sum exchange yields
\[
\begin{aligned}
 \Psi^P_{H,\alpha}(b)
 &=\sum_{t<H}\alpha^t
   \sum_x(P^t)(b,x)g(x)\\
 &=\sum_x
   \left[\sum_{t<H}\alpha^t(P^t)(b,x)\right]g(x)\\
 &=\sum_xW^P_{H,\alpha}(b,x)g(x).
\end{aligned}
\]
Outside the advantage region \(g(x)=0\).  Inside it, the assumed occupation
difference and the gap are nonnegative.  Termwise comparison proves
\cref{thm:kernel-occupation-alignment}.

For the exact persistence boundary, \(H=2\) gives
\(\Psi_2=g+\alpha P(\theta)g\).  At initial belief zero and
\(\alpha=1/2\),
\[
\begin{array}{ccl}
g=(1,0)&:&\Psi_2=1+\theta/2,\\
g=(0,1)&:&\Psi_2=(1-\theta)/2,\\
g=(1,1)&:&\Psi_2=3/2.
\end{array}
\]
Substitution of \(\theta=1/4\) and \(3/4\) gives respectively
\(9/8\to11/8\), \(3/8\to1/8\), and \(3/2\to3/2\).  Both matrices are exactly
row stochastic.  The sign changes because persistence reallocates occupation;
it does not order occupation on the positive-gap region without an alignment
condition.
